\documentclass[12pt]{article}
\usepackage[utf8]{inputenc}
\usepackage[T1]{fontenc}
\usepackage{hyperref}
\title{Tagbanwa Language Resources for K-12 Curriculum Development}
\author{Compiled by Claude Sonnet 5 (claude-sonnet-5) for Andr\'as Kornai}
\date{}
\begin{document}
\maketitle
\section{Introduction}
ISO 639-3 tbw is listed by Wikipedia as Aborlan Tagbanwa, an Austronesian language of Palawan, Philippines. It should not be confused with Central Tagbanwa (tgt) or Calamian Tagbanwa. Sources describe Tagbanwa languages as losing younger speakers to Cuyonon and Tagalog. Little language-specific material was found.
\section{Primary Educational Resources}
No primers or textbooks were found. Wikipedia mentions proposals to teach the Tagbanwa script in schools with Tagbanwa populations; I saw no confirmed programme.
\section{Linguistic and Reference Materials}
\begin{itemize}
\item Wikipedia, Aborlan Tagbanwa (ISO 639:tbw). \url{https://wikipedia.com/wiki/ISO_639:tbw}
\item Tagbanwa script (abugida; Unicode block U+1760--U+177F). \url{https://en.wikipedia.org/wiki/Tagbanwa_script}
\item Central Tagbanwa language (a different language, tgt, for comparison). \url{https://en.wikipedia.org/wiki/Central_Tagbanwa_language}
\item Scebold (2003), \emph{Central Tagbanwa: A Philippine Language on the Brink of Extinction}, concerns Central Tagbanwa, not tbw; no URL was seen.
\end{itemize}
\section{Religious and Community Resources}
None found.
\section{Digital and Audio Resources}
None found. A Tagbanwa-language dictionary page at tagalognow.com appeared in results, but I did not verify its content or quality: \url{https://www.tagalognow.com/dictionary/Tagbanwa/}
\section{Curriculum Development Context}
The script is encoded in Unicode and is a cultural asset for a curriculum. Beyond that, documentation specific to Aborlan Tagbanwa was not found; the curriculum would need community-led documentation and verification of which Tagbanwa variety is meant.
\end{document}